% 对应单元：QM-C04-S01
% 主要来源：Shankar 练习 4.2.1，印刷页 p. 129，PDF p. 141
% 人工核验：是

\clearpage
\section{第 4.1 节代表性教材习题与订正}
\label{exercise-set:QM-C04-S01}

\exercisemeta
  {QM-C04-S01-EX}
  {2026-09-08}
  {Shankar，Exercise 4.2.1 (Very Important)；印刷页 p.~129；PDF p.~141}
  {完成教材习题的状态、谱、测量与相位计算}
  {全部答案与概率归一化已核对}

\exercisequestion{QM-C04-S01-E01}{教材练习 4.2.1：三维空间中的测量与相对相位}

\textbf{对应知识点：}
\begin{itemize}[leftmargin=*]
  \item \relatedtopic{QM-C04-S01-T01}{状态、整体相位与相对相位}；
  \item \relatedtopic{QM-C04-S01-T02}{Hermitian 算符及本征问题}；
  \item \relatedtopic{QM-C04-S01-T03}{Born 规则、简并投影与测量更新}。
\end{itemize}

\textbf{来源说明。}第 4.1 节只集中陈述四条公设，没有独立习题。本题紧邻本节，教材将其标为 “Very Important”；虽然题号属于第 4.2 节，但只使用已经学过的线性代数和第 4.1 节测量公设，因此作为本节代表题提前完成。

在三维复 Hilbert 空间中，以 $L_z$ 的正交归一本征基
\[
  \{\ket{L_z=1},\ket{L_z=0},\ket{L_z=-1}\}
\]
表示状态，并采用 $\hbar=1$ 的无量纲约定。给定
\begin{align}
  L_x&=\frac1{\sqrt2}
  \begin{pmatrix}
    0&1&0\\
    1&0&1\\
    0&1&0
  \end{pmatrix},
  &
  L_y&=\frac1{\sqrt2}
  \begin{pmatrix}
    0&-\ii&0\\
    \ii&0&-\ii\\
    0&\ii&0
  \end{pmatrix},
  &
  L_z&=
  \begin{pmatrix}
    1&0&0\\
    0&0&0\\
    0&0&-1
  \end{pmatrix}.
  \label{eq:QM-C04-S01-E01-angular-momentum-matrices}
\end{align}

完成下列任务：
\begin{enumerate}[leftmargin=*]
  \item 写出测量 $L_z$ 的所有可能结果。
  \item 对状态 $\ket{L_z=1}$，计算 $\langle L_x\rangle$、$\langle L_z\rangle$ 与
  \[
    \Delta L_x
    =\sqrt{\langle L_x^2\rangle-\langle L_x\rangle^2}.
  \]
  \item 在 $L_z$ 基中求 $L_x$ 的全部归一化本征态与本征值。
  \item 初态为 $\ket{L_z=-1}$。若测量 $L_x$，求所有可能结果及其概率。
  \item 初态为
  \[
    \ket{\psi}
    =\begin{pmatrix}1/2&1/2&1/\sqrt2\end{pmatrix}^{\mathsf T}.
  \]
  若先测量 $L_z^2$ 并得到 $+1$，求该结果的概率和测量后状态；另求直接测量 $L_x$ 时各结果的概率。
  \item 已知某归一化态满足
  \[
    P(L_z=1)=\frac14,
    \quad P(L_z=0)=\frac12,
    \quad P(L_z=-1)=\frac14.
  \]
  写出最一般的纯态，说明哪些相位具有物理意义，并分别计算 $P(L_x=0)$ 与 $P(L_x=1)$ 对相位的依赖。
\end{enumerate}

\begin{checkedsolution}
$L_z$ 已经是对角矩阵，所以可能结果就是它的本征值：
\begin{equation}
  \boxed{L_z=1,\ 0,\ -1.}
  \label{eq:QM-C04-S01-E01-Lz-outcomes}
\end{equation}

对
\[
  \ket{L_z=1}=\begin{pmatrix}1&0&0\end{pmatrix}^{\mathsf T},
\]
直接得到
\begin{equation}
  \langle L_x\rangle=0,
  \qquad
  \langle L_z\rangle=1.
  \label{eq:QM-C04-S01-E01-first-moments}
\end{equation}
又因为
\[
  L_x\ket{L_z=1}
  =\frac1{\sqrt2}\begin{pmatrix}0&1&0\end{pmatrix}^{\mathsf T},
\]
所以
\[
  \langle L_x^2\rangle
  =\left\|L_x\ket{L_z=1}\right\|^2
  =\frac12.
\]
故
\begin{equation}
  \boxed{\Delta L_x=\frac1{\sqrt2}.}
  \label{eq:QM-C04-S01-E01-Lx-uncertainty}
\end{equation}

令 $L_x$ 的本征向量为 $(a,b,c)^{\mathsf T}$。特征方程为
\begin{equation}
  \det(L_x-\lambda I)
  =-\lambda(\lambda^2-1)=0,
  \label{eq:QM-C04-S01-E01-Lx-characteristic}
\end{equation}
因此 $\lambda=1,0,-1$。一组正交归一本征态为
\begin{align}
  \ket{L_x=1}
  &=\frac12
    \begin{pmatrix}1\\\sqrt2\\1\end{pmatrix},
  &
  \ket{L_x=0}
  &=\frac1{\sqrt2}
    \begin{pmatrix}1\\0\\-1\end{pmatrix},
  &
  \ket{L_x=-1}
  &=\frac12
    \begin{pmatrix}1\\-\sqrt2\\1\end{pmatrix}.
  \label{eq:QM-C04-S01-E01-Lx-eigenvectors}
\end{align}
每个本征态都可另乘任意整体相位，不影响下列概率。

把初态 $\ket{L_z=-1}=(0,0,1)^{\mathsf T}$ 展开到 $L_x$ 本征基：
\begin{equation}
  \ket{L_z=-1}
  =\frac12\ket{L_x=1}
   -\frac1{\sqrt2}\ket{L_x=0}
   +\frac12\ket{L_x=-1}.
  \label{eq:QM-C04-S01-E01-zminus-in-x-basis}
\end{equation}
故
\begin{equation}
  \boxed{
  P(L_x=1)=\frac14,
  \quad P(L_x=0)=\frac12,
  \quad P(L_x=-1)=\frac14.}
  \label{eq:QM-C04-S01-E01-zminus-Lx-probabilities}
\end{equation}
展开系数 $1/2,-1/\sqrt2,1/2$ 是概率振幅；$1/4,1/2,1/4$ 已经是概率，不能再平方。中间负号影响状态重构和干涉，但概率本身没有负号。

第五问中
\[
  L_z^2=\operatorname{diag}(1,0,1),
  \qquad
  \Pi_{+1}=\operatorname{diag}(1,0,1).
\]
因此
\begin{align}
  P(L_z^2=1)
  &=\bra{\psi}\Pi_{+1}\ket{\psi}
    =\frac14+\frac12
    =\boxed{\frac34},
  \label{eq:QM-C04-S01-E01-Lz2-probability}\\
  \ket{\psi'}
  &=\frac{\Pi_{+1}\ket{\psi}}{\sqrt{3/4}}
    =\boxed{
    \begin{pmatrix}
      1/\sqrt3\\0\\\sqrt{2/3}
    \end{pmatrix}}.
  \label{eq:QM-C04-S01-E01-Lz2-poststate}
\end{align}
这里本征值 $+1$ 二重简并，测量只把状态投影到由 $\ket{L_z=1}$ 与 $\ket{L_z=-1}$ 张成的子空间，不会凭空选定其中某一个方向。

若直接测量 $L_x$，将原初态投影到式~\eqref{eq:QM-C04-S01-E01-Lx-eigenvectors}，得到
\begin{align}
  \braket{L_x=1}{\psi}
  &=\frac{1+2\sqrt2}{4},
  &
  \braket{L_x=0}{\psi}
  &=\frac{1-\sqrt2}{2\sqrt2},
  &
  \braket{L_x=-1}{\psi}
  &=\frac14.
  \label{eq:QM-C04-S01-E01-general-state-Lx-amplitudes}
\end{align}
所以
\begin{equation}
  \boxed{
  P(1)=\frac{9+4\sqrt2}{16},
  \quad
  P(0)=\frac{3-2\sqrt2}{8},
  \quad
  P(-1)=\frac1{16}.}
  \label{eq:QM-C04-S01-E01-general-state-Lx-probabilities}
\end{equation}
三者之和为 $1$。

满足第六问 $L_z$ 概率分布的最一般归一化纯态为
\begin{equation}
  \boxed{
  \ket{\psi}
  =\frac{\ee^{\ii\delta_1}}2\ket{L_z=1}
   +\frac{\ee^{\ii\delta_2}}{\sqrt2}\ket{L_z=0}
   +\frac{\ee^{\ii\delta_3}}2\ket{L_z=-1}.}
  \label{eq:QM-C04-S01-E01-phase-state}
\end{equation}
共同变换 $\delta_j\mapsto\delta_j+\chi$ 只给整个 ket 增加整体相位，故三个相位中只有两个独立的相对相位可能被观测。

对 $L_x=0$，
\begin{align}
  \braket{L_x=0}{\psi}
  &=\frac{\ee^{\ii\delta_1}-\ee^{\ii\delta_3}}{2\sqrt2},
  \notag\\
  P(L_x=0)
  &=\boxed{
  \frac14\left[1-\cos(\delta_1-\delta_3)\right]}
   =\frac12\sin^2\frac{\delta_1-\delta_3}{2}.
  \label{eq:QM-C04-S01-E01-Lx0-phase-probability}
\end{align}
因为 $\ket{L_x=0}$ 在 $\ket{L_z=0}$ 方向上没有分量，此概率与 $\delta_2$ 无关。

对学习对话中进一步检查的 $L_x=1$，
\begin{align}
  \braket{L_x=1}{\psi}
  &=\frac14\left(
    \ee^{\ii\delta_1}
    +2\ee^{\ii\delta_2}
    +\ee^{\ii\delta_3}
  \right),
  \label{eq:QM-C04-S01-E01-Lx1-phase-amplitude}\\
  P(L_x=1)
  &=\boxed{\frac1{16}\left[
    6
    +4\cos(\delta_1-\delta_2)
    +4\cos(\delta_2-\delta_3)
    +2\cos(\delta_1-\delta_3)
  \right]}.
  \label{eq:QM-C04-S01-E01-Lx1-phase-probability}
\end{align}
这里三个两两相位关系都会出现，但它们仍是相位\textbf{差}而不是相位和。原因是 Born 规则使用复数模平方：
\begin{equation}
  \abs{A}^2=A^*A,
  \qquad
  \ee^{-\ii\delta_j}\ee^{\ii\delta_k}
  =\ee^{\ii(\delta_k-\delta_j)}.
  \label{eq:QM-C04-S01-E01-modulus-square-phase-difference}
\end{equation}
普通平方 $A^2$ 的确会产生 $\delta_j+\delta_k$，但它一般是复数，不能作为概率。整体相位自检也给出同一结论：$\delta_j\mapsto\delta_j+\chi$ 不应改变任何概率，而相位差在该变换下保持不变。
\end{checkedsolution}

\textbf{检查点：}先区分本征值、概率振幅、概率和测量后状态。概率必须对振幅取模平方 $A^*A$，不能取普通平方 $A^2$。遇到简并结果时投影到完整本征子空间；判断相位是否可观测时，先去掉共同整体相位，再检查测量基是否把相对相位转化为交叉项。

\clearpage
\section{第 4.2 节教材习题与订正}
\label{exercise-set:QM-C04-S02}

\exercisemeta
  {QM-C04-S02-EX}
  {2026-09-10}
  {Shankar, Exercises 4.2.2--4.2.3；印刷页 p.~139；PDF p.~151}
  {实波函数的平均动量、定义域条件与位置相位造成的动量平移}
  {两题均已完成位置表象、动量表象与算符形式核对}

\exercisequestion{QM-C04-S02-E01}{教材练习 4.2.2：实波函数的平均动量}

\textbf{对应知识点：}
\relatedtopic{QM-C04-S02-T04}{期望值与不确定度}；
\relatedtopic{QM-C04-S02-T06}{高斯波包与动量表象}。

\textbf{题目。}设一维归一化波函数 $\psi(x)$ 为实函数，并满足使动量期望值存在的可微性、收敛性与边界条件。
\begin{enumerate}[label=(\arabic*)]
  \item 证明 $\langle P\rangle=0$。
  \item 推广到 $\psi(x)=c\psi_r(x)$，其中 $\psi_r$ 为实函数，$c\neq0$ 为任意复常数。
  \item 分别说明位置表象分部积分和动量表象奇偶性证明所需的条件。
\end{enumerate}

\begin{checkedsolution}
在位置表象中 $P=-\ii\hbar\,\dd/\dd x$。若 $\psi$ 为实函数，则
\begin{align}
  \bra{\psi}P\ket{\psi}
  &=-\ii\hbar
    \int_{-\infty}^{\infty}\psi(x)\psi'(x)\dd x
  \notag\\
  &=-\frac{\ii\hbar}{2}
    \left[\psi^2(x)\right]_{-\infty}^{\infty}.
  \label{eq:QM-C04-S02-E01-position-proof}
\end{align}
若边界项为零，便有
\begin{equation}
  \boxed{\langle P\rangle=0}.
  \label{eq:QM-C04-S02-E01-zero-mean-momentum}
\end{equation}
一个干净的充分条件是 $\psi\in H^1(\real)$：$\psi,\psi'\in L^2(\real)$，取适当连续代表后 $\psi(x)\to0$（$\abs{x}\to\infty$），且动量矩阵元有定义。仅有 $\psi\in L^2$ 不足以保证逐点边界极限和导数性质。

若 $\psi=c\psi_r$ 且 $\psi_r$ 为实函数，则
\begin{equation}
  \bra{\psi}P\ket{\psi}
  =-\frac{\ii\hbar\abs{c}^2}{2}
   \left[\psi_r^2(x)\right]_{-\infty}^{\infty}=0.
  \label{eq:QM-C04-S02-E01-complex-constant}
\end{equation}
$c$ 任意时 $\psi$ 未必归一化，严格期望值应写为
\begin{equation}
  \langle P\rangle_\psi
  =\frac{\bra{\psi}P\ket{\psi}}{\braket{\psi}{\psi}}=0.
  \label{eq:QM-C04-S02-E01-unnormalized-expectation}
\end{equation}

教材提示给出等价的动量表象证明。对实 $\psi(x)$，Fourier transform 满足
\begin{equation}
  \widetilde\psi(-p)=\widetilde\psi(p)^*,
  \qquad
  \abs{\widetilde\psi(-p)}^2
  =\abs{\widetilde\psi(p)}^2.
  \label{eq:QM-C04-S02-E01-fourier-symmetry}
\end{equation}
若 $\int\abs{p}\abs{\widetilde\psi(p)}^2\dd p<\infty$，则动量概率密度为偶函数，而 $p$ 为奇函数，故
\begin{equation}
  \langle P\rangle
  =\int_{-\infty}^{\infty}
    p\abs{\widetilde\psi(p)}^2\dd p=0.
  \label{eq:QM-C04-S02-E01-momentum-proof}
\end{equation}
\end{checkedsolution}

\exercisequestion{QM-C04-S02-E02}{教材练习 4.2.3：位置相位与动量平移}

\textbf{对应知识点：}
\relatedtopic{QM-C04-S02-T06}{Fourier 变换与动量分布}，以及位置相关的相对相位。

\textbf{题目。}设归一化波函数 $\psi(x)$ 的平均动量为 $\bar p$。定义
\begin{equation}
  \psi_{p_0}(x)=\ee^{\ii p_0x/\hbar}\psi(x),
  \qquad p_0\in\real.
  \label{eq:QM-C04-S02-E02-boosted-wavefunction}
\end{equation}
证明新波函数仍归一化、平均动量为 $\bar p+p_0$，并说明其位置和动量概率分布怎样变化。

\begin{checkedsolution}
由于 $\abs{\ee^{\ii p_0x/\hbar}}=1$，
\begin{equation}
  \int_{-\infty}^{\infty}\abs{\psi_{p_0}(x)}^2\dd x
  =\int_{-\infty}^{\infty}\abs{\psi(x)}^2\dd x=1.
  \label{eq:QM-C04-S02-E02-normalization}
\end{equation}
同一等式说明位置概率密度完全不变。该因子依赖 $x$，不是可忽略的整体相位；它会被动量导数算符读出：
\begin{align}
  P\psi_{p_0}
  &=-\ii\hbar\frac{\dd}{\dd x}
    \left(\ee^{\ii p_0x/\hbar}\psi\right)
  \notag\\
  &=\ee^{\ii p_0x/\hbar}(p_0+P)\psi.
  \label{eq:QM-C04-S02-E02-momentum-action}
\end{align}
因此
\begin{equation}
  \boxed{\langle P\rangle_{p_0}=p_0+\bar p}.
  \label{eq:QM-C04-S02-E02-shifted-mean}
\end{equation}
若 $\psi\in H^1(\real)$，有限 $p_0$ 下乘以此光滑相位仍属于 $H^1(\real)$，故动量算符定义域保持。

在动量表象中，
\begin{align}
  \widetilde\psi_{p_0}(p)
  &=\frac1{\sqrt{2\pi\hbar}}
    \int\ee^{-\ii px/\hbar}
      \ee^{\ii p_0x/\hbar}\psi(x)\dd x
  \notag\\
  &=\boxed{\widetilde\psi(p-p_0)}.
  \label{eq:QM-C04-S02-E02-fourier-shift}
\end{align}
故 $\abs{\widetilde\psi(p)}^2$ 整体向 $+p$ 方向平移 $p_0$，形状和动量不确定度不变。

令 $U(p_0)=\ee^{\ii p_0X/\hbar}$，同一结论可压缩为
\begin{equation}
  \boxed{U^\dagger(p_0)PU(p_0)=P+p_0I},
  \label{eq:QM-C04-S02-E02-operator-identity}
\end{equation}
因此 $U(p_0)$ 是 momentum translation（动量平移）或 momentum boost 算符。
\end{checkedsolution}

\textbf{检查点：}实波函数给出关于 $p=0$ 对称的动量分布，但位置表象证明必须说明动量算符定义域和边界项。乘 $\ee^{\ii p_0x/\hbar}$ 不改变位置概率密度，却会把动量分布整体平移；判断相位是否可观测，关键在于它是否随位置变化并被相应算符读取。

\clearpage
\section{第 4.3 节正文推导、应用与综合检验}
\label{exercise-set:QM-C04-S03}

\exercisemeta
  {QM-C04-S03-EX}
  {2026-09-11}
  {Shankar 第 4.3 节正文；印刷页 pp.~145--150；PDF pp.~157--162；另含一项自编综合检验}
  {传播算符、恒力势的动量表象、分段含时 Hamiltonian}
  {两项教材正文推导或应用与一项自编综合检验均已完成并订正}

\textbf{来源说明。}本节教材没有独立编号的例题或习题。以下 E01 基于教材式 (4.3.14)--(4.3.16) 的正文推导，E02 是教材末尾明确提出的恒力势应用；E03 为用于检查时间排序的自编题，不是教材原题。

\exercisequestion{QM-C04-S03-E01}{教材正文推导：传播算符的基本性质}

\textbf{对应知识点：}
\relatedtopic{QM-C04-S03-T02}{时间无关 Hamiltonian、定态与传播算符}。

\textbf{题目。}设 $H$ 为时间无关 Hermitian 算符，定义
\begin{equation}
  U(t_2,t_1)=\exp\left[-\frac{\ii}{\hbar}H(t_2-t_1)\right].
  \label{eq:QM-C04-S03-E01-definition}
\end{equation}
验证它满足 Schrödinger 方程和初始条件，并证明传播算符的复合律、逆算符关系以及内积保持性。

\begin{checkedsolution}
对 $t_2$ 求导。由于指数中只有同一个 $H$，$H$ 与其指数函数对易，故
\begin{align}
  \ii\hbar\frac{\partial}{\partial t_2}U(t_2,t_1)
  &=\ii\hbar\left(-\frac{\ii}{\hbar}H\right)
    U(t_2,t_1)\notag\\
  &=HU(t_2,t_1).
  \label{eq:QM-C04-S03-E01-evolution-equation}
\end{align}
当 $t_2=t_1$ 时，
\begin{equation}
  U(t_1,t_1)=\ee^0=I.
  \label{eq:QM-C04-S03-E01-initial-condition}
\end{equation}

若 $t_3\ge t_2\ge t_1$，两个指数仍是同一算符 $H$ 的函数，故可合并：
\begin{align}
  U(t_3,t_2)U(t_2,t_1)
  &=\ee^{-\ii H(t_3-t_2)/\hbar}
    \ee^{-\ii H(t_2-t_1)/\hbar}\notag\\
  &=\ee^{-\ii H(t_3-t_1)/\hbar}
   =\boxed{U(t_3,t_1)}.
  \label{eq:QM-C04-S03-E01-composition}
\end{align}
取伴随并使用 $H^\dagger=H$，
\begin{equation}
  U^\dagger(t_2,t_1)
  =\ee^{\ii H(t_2-t_1)/\hbar}
  =U(t_1,t_2)=U^{-1}(t_2,t_1).
  \label{eq:QM-C04-S03-E01-adjoint-inverse}
\end{equation}
所以 $U^\dagger U=I$。对任意两个同时演化的态，
\begin{align}
  \braket{\phi(t_2)}{\psi(t_2)}
  &=\bra{\phi(t_1)}U^\dagger(t_2,t_1)
    U(t_2,t_1)\ket{\psi(t_1)}\notag\\
  &=\braket{\phi(t_1)}{\psi(t_1)}.
  \label{eq:QM-C04-S03-E01-inner-product}
\end{align}
因此范数、总概率和任意两态间的内积均保持。
\end{checkedsolution}

\exercisequestion{QM-C04-S03-E02}{教材正文应用：恒力势的动量表象本征函数}

\textbf{对应知识点：}
\relatedtopic{QM-C04-S03-T04}{表象选择与构型空间}。

\textbf{题目。}对一维恒力势
\begin{equation}
  H=\frac{P^2}{2m}-fX,
  \qquad f>0,
  \label{eq:QM-C04-S03-E02-hamiltonian}
\end{equation}
在动量表象中求能量本征函数 $\widetilde\psi_E(p)=\braket{p}{E}$，说明它为何不能作普通 $L^2$ 归一化，并确定使其满足 $\braket{E}{E'}=\delta(E-E')$ 的归一化常数。

\begin{checkedsolution}
动量表象中 $P\mapsto p$、$X\mapsto\ii\hbar\,\dd/\dd p$，故
\begin{equation}
  \left(\frac{p^2}{2m}-\ii\hbar f\frac{\dd}{\dd p}\right)
  \widetilde\psi_E(p)=E\widetilde\psi_E(p).
  \label{eq:QM-C04-S03-E02-eigen-equation}
\end{equation}
整理为可分离的一阶方程：
\begin{equation}
  \frac{1}{\widetilde\psi_E}
  \frac{\dd\widetilde\psi_E}{\dd p}
  =\frac{\ii}{\hbar f}
    \left(E-\frac{p^2}{2m}\right).
  \label{eq:QM-C04-S03-E02-separated-ode}
\end{equation}
两边积分，
\begin{align}
  \ln\widetilde\psi_E(p)
  &=\frac{\ii}{\hbar f}
    \left(Ep-\frac{p^3}{6m}\right)+\ln C,
  \label{eq:QM-C04-S03-E02-log-solution}\\
  \widetilde\psi_E(p)
  &=\boxed{C\exp\left[\frac{\ii}{\hbar f}
    \left(Ep-\frac{p^3}{6m}\right)\right]}.
  \label{eq:QM-C04-S03-E02-solution}
\end{align}
指数是纯相位，所以 $\abs{\widetilde\psi_E(p)}=\abs C$ 在整个实轴恒定，不能使 $\int\abs{\widetilde\psi_E}^2\dd p=1$。它与平面波类似，是连续谱的广义本征函数，应作 Dirac delta 归一化。

计算两个能量态的内积时，$p^3$ 相位相消：
\begin{align}
  \braket{E}{E'}
  &=\abs C^2\int_{-\infty}^{\infty}
    \exp\left[\frac{\ii(E'-E)p}{\hbar f}\right]\dd p\notag\\
  &=\abs C^2(2\pi\hbar f)\delta(E'-E).
  \label{eq:QM-C04-S03-E02-delta-normalization}
\end{align}
因此可取
\begin{equation}
  \boxed{C=\frac1{\sqrt{2\pi\hbar f}}}
  \label{eq:QM-C04-S03-E02-normalization-constant}
\end{equation}
另乘任意只依赖 $E$ 的整体相位不影响归一化。若允许 $f<0$，则应改为
\begin{equation}
  \abs C=\frac1{\sqrt{2\pi\hbar\abs f}}.
  \label{eq:QM-C04-S03-E02-normalization-general-force}
\end{equation}
\end{checkedsolution}

\exercisequestion{QM-C04-S03-E03}{自编综合检验：分段二能级 Hamiltonian}

\textbf{对应知识点：}
\relatedtopic{QM-C04-S03-T03}{含时 Hamiltonian 与时间排序}。

\textbf{题目。}在基 $\{\ket{1},\ket{2}\}$ 中令
\begin{equation}
  Z=\begin{pmatrix}1&0\\0&-1\end{pmatrix},
  \qquad
  X=\begin{pmatrix}0&1\\1&0\end{pmatrix}.
  \label{eq:QM-C04-S03-E03-pauli-matrices}
\end{equation}
系统初态为 $\ket{\psi(0)}=\ket{1}$，Hamiltonian 分段取值
\begin{equation}
  H(t)=
  \begin{cases}
    \hbar\omega Z,&0\le t<T,\\
    \hbar\omega X,&T\le t<2T,
  \end{cases}
  \qquad \omega T=\frac{\pi}{4}.
  \label{eq:QM-C04-S03-E03-piecewise-hamiltonian}
\end{equation}
求 $\ket{\psi(2T)}$、在 $\ket{1},\ket{2}$ 基中的测量概率及在 $\ket{+}=(\ket{1}+\ket{2})/\sqrt2$ 上的概率；再交换两个时间段的顺序，比较结果。

\begin{checkedsolution}
利用 $Z^2=X^2=I$，不能把矩阵指数理解为对矩阵元逐项取指数；应使用
\begin{equation}
  \ee^{-\ii\theta A}=I\cos\theta-\ii A\sin\theta,
  \qquad A^2=I.
  \label{eq:QM-C04-S03-E03-exponential-identity}
\end{equation}
因此
\begin{align}
  U_1&=\ee^{-\ii\pi Z/4}
  =\frac1{\sqrt2}(I-\ii Z)
  =\frac1{\sqrt2}
    \begin{pmatrix}1-\ii&0\\0&1+\ii\end{pmatrix},
  \label{eq:QM-C04-S03-E03-U1}\\
  U_2&=\ee^{-\ii\pi X/4}
  =\frac1{\sqrt2}(I-\ii X)
  =\frac1{\sqrt2}
    \begin{pmatrix}1&-\ii\\-\ii&1\end{pmatrix}.
  \label{eq:QM-C04-S03-E03-U2}
\end{align}
较晚的第二段位于左侧，所以正确总传播算符是 $U_2U_1$：
\begin{align}
  \ket{\psi(2T)}
  &=U_2U_1\ket{1}
   =\boxed{\frac12
    \begin{pmatrix}1-\ii\\-1-\ii\end{pmatrix}}
  \notag\\
  &=\frac{\ee^{-\ii\pi/4}}{\sqrt2}
    \begin{pmatrix}1\\-\ii\end{pmatrix}.
  \label{eq:QM-C04-S03-E03-correct-state}
\end{align}
故
\begin{equation}
  P(1)=P(2)=\frac12.
  \label{eq:QM-C04-S03-E03-z-probabilities}
\end{equation}
在 $\ket{+}$ 上的振幅与概率为
\begin{align}
  \braket{+}{\psi(2T)}
  &=\frac{(1-\ii)+(-1-\ii)}{2\sqrt2}
   =-\frac{\ii}{\sqrt2},
  \label{eq:QM-C04-S03-E03-plus-amplitude}\\
  P(+)&=\boxed{\frac12}.
  \label{eq:QM-C04-S03-E03-plus-probability}
\end{align}

若交换两个时间段，传播顺序变为 $U_1U_2$：
\begin{equation}
  U_1U_2\ket{1}
  =\frac{\ee^{-\ii\pi/4}}{\sqrt2}
    \begin{pmatrix}1\\1\end{pmatrix}
  =\ee^{-\ii\pi/4}\ket{+}.
  \label{eq:QM-C04-S03-E03-reversed-state}
\end{equation}
此时仍有 $P(1)=P(2)=1/2$，但
\begin{equation}
  \boxed{P_{\mathrm{reversed}}(+)=1}.
  \label{eq:QM-C04-S03-E03-reversed-plus-probability}
\end{equation}
两种顺序在 $\ket{1},\ket{2}$ 测量中恰好给出相同概率，却在另一测量基中可区分；这说明只比较一组概率不足以判定量子态相同，也直接体现了 $[X,Z]\neq0$ 时演化次序的重要性。
\end{checkedsolution}

\textbf{检查点：}时间无关传播算符的复合来自同一 $H$ 的指数可合并；含时情形则必须保持时间次序。连续谱本征函数通常不是 $L^2$ 态，应使用 delta 归一化。计算矩阵指数时先利用谱分解或算符满足的多项式关系，不能逐元素取指数。
